\documentclass[11pt,reqno]{amsart}
\usepackage[T1]{fontenc}
\usepackage[utf8]{inputenc}
\usepackage{amsmath,amssymb,amsthm}
\usepackage[margin=1.15in]{geometry}
\usepackage{booktabs}
\usepackage{array}
\usepackage[hidelinks]{hyperref}

\theoremstyle{plain}
\newtheorem{theorem}{Theorem}[section]
\newtheorem{proposition}[theorem]{Proposition}
\newtheorem{lemma}[theorem]{Lemma}
\newtheorem{corollary}[theorem]{Corollary}
\newtheorem{conjecture}[theorem]{Conjecture}
\newtheorem{criterion}[theorem]{Criterion}
\theoremstyle{definition}
\newtheorem{definition}[theorem]{Definition}
\newtheorem{remark}[theorem]{Remark}
\newtheorem{observation}[theorem]{Observation}
\newtheorem{formulation}[theorem]{Formulation}

\newcommand{\Zt}{\mathbb{Z}_2}
\newcommand{\Z}{\mathbb{Z}}
\newcommand{\Q}{\mathbb{Q}}
\newcommand{\R}{\mathbb{R}}
\newcommand{\N}{\mathbb{N}}
\newcommand{\vt}{\nu_2}
\DeclareMathOperator{\Ph}{\Phi}
\newcommand{\lab}[1]{\textnormal{\textsc{\small #1}}}

\begin{document}

\title[After the reduction]
{After the reduction: proved statistics, excluded structure,\\
and the missing pointwise tool in the accelerated $3x+1$ problem}

\author{Elias De Jes\'us}
\address{Independent Researcher}
\email{dejesuselias10@gmail.com}
\urladdr{https://orcid.org/0009-0007-0190-9143}
\date{September 2026}

\subjclass[2020]{Primary 11B83; Secondary 11J72, 11K31, 37P05, 68R15, 03B35}
\keywords{$3x+1$ problem, accelerated Collatz map, 2-adic integers, confinement,
Sturmian words, formalized mathematics}

\begin{abstract}
For the accelerated $3x+1$ map $T(m)=(3m+1)/2^{\vt(3m+1)}$ on odd integers, the
existence of a divergent orbit is \emph{equivalent} to the existence of a single
positive odd seed whose cumulative valuation satisfies $2^{S_n}\le 3^n$ at every
depth --- a statement that is machine-checked and carries no hypotheses. This note
reports what the resulting programme has established and, more usefully, what it
has excluded. At the population level the statistics are complete: the confined
words have a proved exponential rate, a proved two-sided polynomial correction,
and the record-holding integers sit where random placement at that rate predicts,
to within a bounded and measured residual. At the structured level the natural
candidates are gone: periodic confined points are negative rationals, the critical
Sturmian word has an irrational conjugacy value, words of non-critical density are
excluded by density alone, and descent at the $3$-adic place cannot accumulate a
budget. What remains is a \emph{placement} problem about specific integers inside
a population whose size is known, and we argue that no tool currently in the
programme can address it: we give a proved obstruction showing that the aggregate
identity, the confinement condition and the sign of a denominator cannot together
separate the positive integers from the rest, because the $3x-1$ system agrees on
all three and disagrees on the answer. We state the three conditions any candidate
tool must satisfy, the controls it must fail, and seven open problems.
\textbf{Nothing here bears on the Collatz conjecture, and no claim of progress
toward it is made or implied.}
\end{abstract}

\maketitle

\section{Introduction}\label{sec:intro}

Write $T$ for the \emph{accelerated} $3x+1$ map on the odd positive integers,
\[
T(m)=\frac{3m+1}{2^{\,a(m)}},\qquad a(m)=\vt(3m+1)\ge 1,
\]
and put $m_k=T^k(m_0)$, $a_k=a(m_k)$, $S_n=\sum_{k<n}a_k$, and
\[
\alpha=\log_2 3=1.58496\ldots,\qquad R_n=S_n-n\alpha .
\]
$R_n$ is the \emph{drift}: by the aggregate identity of Section~\ref{sec:boundary},
$m_n$ is comparable to $m_0\,2^{-R_n}$, so an orbit grows exactly when its drift
falls. The whole subject sits on the line $R_n=0$.

\medskip
\noindent\textbf{The thesis, in one paragraph.} The divergence question reduces,
exactly and unconditionally, to the existence of one positive odd integer whose
drift never rises above zero (Theorem~\ref{thm:reduction}). After that reduction
the \emph{statistics} of such confinement are settled --- an exact exponential
rate, an exact polynomial correction, and record-holding integers that sit where
random placement at that rate predicts --- and the \emph{structured} candidates
are excluded one family at a time: periodic, algebraic, Sturmian, bounded-drift,
and descent-accessible. What is left is neither a counting problem nor a
structural one. It is a question about where particular integers lie inside a
population whose size is already known, and the tools that decided every other
part of the programme are provably unable to see it. Section~\ref{sec:needed} is
an attempt to say precisely what a tool that could see it would have to do.

\medskip
\noindent\textbf{Explicit non-claims.} This note does not prove, and does not
claim progress toward, the Collatz conjecture; nor does it exclude divergent
orbits, nor nontrivial cycles. The central reduction is an \emph{equivalence},
and Section~\ref{sec:closed} records a list of reformulations that were mistaken
for progress until they were proved equivalent to the problem itself. Where a
statement rests on an unrefereed source this is said at every occurrence, not
once.

\medskip
\noindent\textbf{Labels.} Every substantive claim carries one of:
\lab{proved (lean)}, machine-checked in the accompanying repository with no
hypotheses and no axioms beyond Lean's three standard ones;
\lab{proved (paper)}, proved by the author on paper and audited, not formalized;
\lab{cited}, an external result, with reference;
\lab{verified}, an exact computation over a stated finite range;
\lab{heuristic}; \lab{conjecture}; \lab{open}. Section~\ref{sec:resources} says
where each machine-checked statement lives and how each computation is
reproduced.

\section{The valuation boundary}\label{sec:boundary}

Everything below rests on one identity, which we display once because it is used
at every later step. Since $2^{a_k}m_{k+1}=3m_k+1$ for each $k$, induction gives
the \emph{aggregate identity}
\begin{equation}\label{eq:aggregate}
2^{\,S_n}\,m_n\;=\;3^{\,n}m_0+C_n,\qquad
C_n=\sum_{i<n}3^{\,n-1-i}\,2^{\,S_i}\;\ge\;3^{\,n-1}>0\quad(n\ge1).
\end{equation}
\lab{proved (lean)} over $\N$, \texttt{Divergence.aggregate\_identity}; the same
induction over $\Zt^{\times}$, which is what Theorem~\ref{thm:marker} needs, is
\lab{proved (paper)}. Two consequences are used repeatedly. Dividing by $2^{S_n}$
and taking logarithms gives Proposition~\ref{prop:drift}. And evaluated at a
period, \eqref{eq:aggregate} becomes a formula for the periodic point, which is
Theorem~\ref{thm:marker}; the positivity of $C_n$, visible in
\eqref{eq:aggregate}, is what fixes its sign.

\begin{proposition}[drift identity]\label{prop:drift}
For odd $m_0>1$, with $E_N=\sum_{k<N}\log_2\!\bigl(1+1/(3m_k)\bigr)$,
\[
\log_2\frac{m_N}{m_0}=N\log_2 3+E_N-S_N,
\qquad 0<E_N\le N\log_2\tfrac43 .
\]
\end{proposition}
\lab{proved (paper)} \cite{DeJesusVM}. The error term $E_N$ is positive and
$O(N)$ with a small constant; for a divergent orbit it is in fact
$O(1)$, since $\sum 1/m_k<\infty$ (Theorem~\ref{thm:summable}). So the sign of
$\log_2(m_N/m_0)$ is decided by $-R_N$ up to a bounded error.

\begin{theorem}[valuation-mean classification]\label{thm:vm}
Let $M$ be a bound below which the $3x+1$ conjecture has been verified and put
$\tau_M=\log_2(3+1/M)$. Every orbit of an odd $m_0>1$ falls into exactly one of
three cases, with disjoint asymptotic valuation means
$\overline D_N=S_N/N$:
\[
\text{terminating: }\lim \overline D_N=2;\qquad
\text{cyclic: }\lim \overline D_N=S/L\in(\log_2 3,\tau_M];\qquad
\text{injective: }\limsup \overline D_N\le\log_2 3 .
\]
Equivalently, $\limsup_N \overline D_N>\tau_M$ if and only if the orbit reaches $1$.
\end{theorem}
\lab{proved (paper)} \cite{DeJesusVM}. The ranges are disjoint because
$\log_2 3<\tau_M<2$. The statement is conditional only through $M$, and only in
the placement of $\tau_M$: the boundary $\log_2 3$ separating injective from
cyclic behaviour is exact and unconditional.

\begin{remark}[why one argument cannot serve both halves]\label{rem:twohalves}
The two open halves of the $3x+1$ problem lie on \emph{opposite sides} of a single
line. A nontrivial cycle would have valuation mean strictly \emph{above}
$\log_2 3$; a divergent orbit has $\limsup$ at or \emph{below} it. Any argument
phrased purely as a bound on the time-average of $a_k$ therefore pushes the two
cases in opposite directions, and cannot exclude both. This is why the programme
splits at the outset into a divergence side and a cycle side, and why the two have
required different machinery throughout.
\end{remark}

\section{The exact reduction}\label{sec:reduction}

\begin{definition}\label{def:confined}
An odd $m$ is \emph{zero-confined} if $2^{S_n(m)}\le 3^n$ for every $n\ge 1$;
equivalently $R_n\le 0$ for every $n$. We write $K$ for the set of zero-confined
elements of $\Zt^{\times}$ --- the \emph{confined dust} --- and $r_{\min}(N,0)$
for the least positive odd integer that is zero-confined for $N$ steps.
\end{definition}

\begin{theorem}[exact reduction]\label{thm:reduction}
There is an odd $M\in\N$ whose accelerated orbit diverges if and only if there is
an odd $m\in\N$ that is zero-confined.
\end{theorem}
\lab{proved (lean)}, \texttt{Divergence.divergent\_iff\_zeroConfined}
\cite{EOCDivergenceRepo}: no hypotheses, axioms \texttt{propext},
\texttt{Classical.choice}, \texttt{Quot.sound}. The forward direction is the
last-maximum argument; the converse is the observation that a zero-confined seed
cannot repeat a value (Theorem~\ref{thm:cycledrift}) and so has an injective,
hence divergent, orbit. The reduction is unconditional because the one external
input it needs --- the windowed collision bound of
Theorem~\ref{thm:sparsity} --- is proved in the same development rather than
assumed.

\begin{definition}[(DE)]\label{def:de}
\emph{Universal drift exit} is the assertion that the right-hand side of
Theorem~\ref{thm:reduction} is false: no positive odd integer is zero-confined.
\lab{open}.
\end{definition}

\begin{proposition}[equivalent forms]\label{prop:equivalent}
The following are equivalent to \textup{(DE)}:
\textup{(a)} the statement of Definition~\ref{def:de};
\textup{(b)} $r_{\min}(N,0)\to\infty$;
\textup{(c)} for some, equivalently every, fixed corridor width $c\ge 0$, no
positive odd integer satisfies $2^{S_n}\le 2^{c}3^{n}$ at every $n$;
\textup{(d)} no zero-confined $2$-adic integer has an expansion ending in
$0^{\infty}$.
\end{proposition}
\lab{proved (lean)} for (a); \lab{proved (paper)} for the equivalence of (b),
(c), (d) \cite{DeJesusEOC,EOCDivergenceRepo}. These are recorded precisely so that
a reformulation into one of them is not mistaken for progress; the repository
keeps the list for that purpose.

\begin{theorem}[the signed marker]\label{thm:marker}
Every eventually periodic point of $K$ is a \emph{negative} rational. Explicitly,
a purely periodic zero-confined point of period $L$ and total $S=S_L$ satisfies
\[
m=\frac{C_L}{2^{S}-3^{L}},\qquad
C_L=\sum_{i<L}3^{\,L-1-i}2^{\,S_i}\ \ge\ 3^{\,L-1}>0,\qquad 2^{S}-3^{L}<0 .
\]
\end{theorem}
\lab{proved (paper)}. Periodicity turns the aggregate identity into
$m(2^{S}-3^{L})=C_L$; confinement at the single index $L$ gives $2^{S}\le 3^{L}$,
strictly by unique factorisation. \lab{verified} on all $1\,717$ zero-confined
purely periodic points of period $\le 10$: every one is negative, without
exception. The marker classifies the rational points of $K$ that can be
enumerated, and says nothing about the rest; its real use is in
Section~\ref{sec:needed}, where the same computation run on $3x-1$ yields the
opposite sign.

\section{The statistics}\label{sec:statistics}

This section is the part of the programme that is finished. Its conclusion is
that at the population level, and at the extreme order statistic, the integers
behave like random placement at the proved rate.

\begin{theorem}[windowed sparsity]\label{thm:sparsity}
There is $\theta$ with $1\le\theta<2$ such that every collision-free
$A\subseteq\N$ satisfies
$\#\bigl(A\cap[a,a+2^{N})\bigr)\le 4(N+1)\theta^{N}$ for all $N$ and all $a\ge1$.
\end{theorem}
\lab{proved (lean)}, \texttt{Divergence.exists\_window\_sparsity}, with
$\theta\approx1.9762$. This is the Fundamental Lemma of Garc\'ia and Tal
\cite{GarciaTal1999} in the explicit, Heppner-free form of Curry
\cite{Curry2026} (\lab{cited} for the originals). Formalizing it is what makes
Theorem~\ref{thm:reduction} unconditional.

\begin{theorem}[reciprocal summability]\label{thm:summable}
A divergent accelerated orbit satisfies $\sum_n 1/m_n<\infty$.
\end{theorem}
\lab{proved (lean)}, \texttt{Divergence.summable\_inv\_orbit}; after
\cite{Curry2026}. It is this that makes $E_N$ in Proposition~\ref{prop:drift}
bounded along a divergent orbit.

\begin{theorem}[cycle drift]\label{thm:cycledrift}
If $m_i=m_j$ with $i<j$ then $3^{\,j-i}<2^{\,S(m_i,\,j-i)}$. In particular no
periodic orbit is zero-confined.
\end{theorem}
\lab{proved (lean)}, \texttt{Divergence.cycle\_drift}. Compare
Theorem~\ref{thm:vm}: a cycle must have valuation mean strictly above $\log_2 3$,
which is the same fact in the other coordinate.

\begin{theorem}[confined-mass rate]\label{thm:rate}
Let $p_N(c)=\sum 2^{-S_N(d)}$, summed over the $c$-confined valuation words $d$ of
length $N$. Then
\[
-\frac1N\log_2 p_N(c)\ \longrightarrow\ I_0=\alpha\bigl(1-H_2(1/\alpha)\bigr)
=0.0793186\ldots
\]
\end{theorem}
\lab{proved (lean)}, \texttt{Occupation.confined\_mass\_rate}. The proof is purely
combinatorial: the upper bound organises words by total valuation and uses
unimodality of $\binom{s-1}{N-1}2^{-s}$; the lower bound is the cycle lemma of
Dvoretzky and Motzkin \cite{DvoretzkyMotzkin1947} (\lab{cited}) applied at the
integral anchor $s_N=\lfloor N\alpha\rfloor$. No probability, no central limit
theorem, no tilting.

\begin{theorem}[three scales]\label{thm:scales}
The polynomial correction to the rate of Theorem~\ref{thm:rate} is exactly
$N^{-3/2}$, two-sided, for both the unweighted count of confined words and the
dyadic survival probability $p_N(c)$; and the $O(1)$ phase multiplying it is an
explicit geometrically weighted renewal sum, whose jump sets are the Beatty
sequences of $\beta=\alpha-1$.
\end{theorem}
\lab{proved (paper)} \cite{DeJesusScales} (version v3, the second revision); not
formalized. Its two-sided local ballot input is derived there from classical
fluctuation theory (\lab{cited}: Caravenna \cite{Caravenna2005}, with the ladder
identity attributed there to Alili and Doney). That note makes \emph{no priority
claim} for the renewal-phase result relative to the announced but unlocated work
of Banderier and Wallner \cite{BanderierWallner}, and neither does this note.

The computation of $p_N(0)$ rests on the prefix correspondence of Terras
\cite{Terras1976} and Everett \cite{Everett1977} (\lab{cited}): the odd $m$
realising a given valuation word of total $S$ form \emph{exactly one} residue
class modulo $2^{S+1}$. Their density among the odd integers is therefore $2^{-S}$
on the nose, which is what makes the following exact rather than asymptotic.

\begin{proposition}[the density, exactly]\label{prop:density}
$p_N(0)$ is computable in integer arithmetic by the transfer recursion
$f[0][0]=1$, $f[k+1][S']=\sum_{d\ge1}f[k][S'-d]$ for $S'\le A[k]$, where
$A[k]=\lfloor k\log_2 3\rfloor$; then
$\mathrm{num}_N=\sum_S f[N][S]\,2^{A[N]-S}$ is an exact integer with
$p_N(0)=\mathrm{num}_N/2^{A[N]}$, and $\mathrm{num}_N$ is \emph{also} the exact
number of odd $m\in[1,2^{A[N]+1})$ that are $N$-confined. Over $N=100,\dots,300$,
\[
\log_2\bigl(1/p_N(0)\bigr)-I_0N-\tfrac32\log_2 N\ =\ -3.28\pm0.05 .
\]
\end{proposition}
\lab{verified}. The identification of $\mathrm{num}_N$ with a count was checked
against direct orbit iteration for $N\le14$ (for instance $\mathrm{num}_2=3$, and
the $2$-confined odd $m<16$ are exactly $\{7,11,15\}$). The displayed residual is
an exact-arithmetic confirmation of the $N^{-3/2}$ factor of
Theorem~\ref{thm:scales} over that range: with the factor the residual is flat to
about a tenth of a bit per doubling, without it the residual grows by about $1.4$
bits per doubling. It is a numerical confirmation, not a proof of the exponent.

\begin{conjecture}[sharp EOC]\label{conj:sharp}
\[
\log_2 r_{\min}(N,0)\ =\ I_0\,N+\tfrac32\log_2 N+O(1).
\]
\end{conjecture}
\lab{conjecture}. Supporting data, \lab{verified}: an exact scan of every odd
$m\equiv3\ (\mathrm{mod}\ 4)$ below $10^{12}$ certifies $r_{\min}(N,0)$ for
$N=1,\dots,346$, with $24$ distinct record holders, the largest being
$r_{\min}(345)=898\,696\,369\,947$. Writing
$\Delta_1(N)=\log_2 r_{\min}(N,0)-\log_2(1/p_N(0))$, the empirical growth exponent
is
\[
I_0+b=0.0737,\qquad\text{block-bootstrap }95\%\text{ CI }[0.061,\ 0.087],
\]
an interval that contains $I_0=0.0793$ and is bounded away from $0$ at both ends.
A persistent drift in $\Delta_1$ and a bounded residual cannot be separated on
this range: the two models differ by $\Delta\mathrm{AICc}=0.29$ at an effective
sample size of about $7$, the record holders being strongly serially correlated.

\begin{proposition}[clustering]\label{prop:clustering}
The confined integers are not a Poisson process at the extreme. Two measurements:
$\Delta_1(N)>0$ at all $24$ record holders, so the \emph{first} confined integer
is $1.9$ to $2.5$ bits further out than random placement predicts; and the direct
count $\#\{\text{odd }m\le 2^{j}:\ N\text{-confined}\}$ exceeds $2^{\,j-1}p_N(0)$
immediately above $r_{\min}(N,0)$ and nowhere else --- by a factor $1.78$ at
$N=200$, $X=2^{28}$, and $5.45$ at $N=300$, $X=2^{35}$ (twelve observed against
$2.2$ expected) --- decaying to $\approx1$ by $2^{40}$, while in the
well-populated regime the null is correct to a few parts in a thousand.
\end{proposition}
\lab{verified} for the two measurements. \lab{heuristic} for the attribution:
these are the two signatures of a clustered point process, and the cause is the
backward chains of Theorem~\ref{thm:descent}, a confined $m\equiv2\ (\mathrm{mod}\ 3)$
having a smaller and more-confined backward image. Thinning the population to
chain roots removes about a quarter to a third of the offset and no more; and it
cannot move $\Delta_1$ at all, because $r_{\min}(N,0)$ is always a chain root
(that is part (L4) of Theorem~\ref{thm:descent}), so the thinned null differs from
the unthinned one by the constant $\log_2(3/2)$ exactly.

\medskip
\noindent\textbf{What Section~\ref{sec:statistics} amounts to.} The size of the
confined population is known exactly at leading order, at polynomial order, and
--- to within a measured constant --- at the extreme. The residual $\Delta_1$ is
positive, bounded over the computable range, and shows no drift that the data can
distinguish from none. Improving any of these estimates does not touch
\textup{(DE)}: a sharper count is still a count.

\section{The structured sectors, excluded}\label{sec:excluded}

The natural way to look for a zero-confined integer is to look for a
\emph{structured} one. This section records that each natural kind of structure
has been excluded, and by what.

\medskip
\noindent\textbf{Algebraic and periodic structure.} By Theorem~\ref{thm:cycledrift}
no periodic orbit is zero-confined, and by Theorem~\ref{thm:marker} every
eventually periodic point of $K$ is a negative rational. So no positive integer
can be confined through periodicity of its orbit.

\medskip
\noindent\textbf{The dictionary.} Two coordinate systems are in play and it is
worth fixing the translation. In \emph{accelerated} coordinates an orbit is
recorded by its valuation word $(a_0,a_1,\dots)$ with $a_k\ge1$. In
\emph{parity} coordinates one uses the standard map
$T_{\mathrm{std}}(x)=x/2$ for even $x$ and $(3x+1)/2$ for odd $x$, and records the
parity word $v_i\equiv T_{\mathrm{std}}^{\,i}(x)\pmod 2$. One accelerated step
consumes one letter $1$ followed by $a_k-1$ letters $0$, so the valuation word is
the sequence of gaps between consecutive ones of the parity word, and
$S_n$ is the position of the $(n+1)$-st one. Writing $k_\ell$ for the number of
ones among $v_0,\dots,v_{\ell-1}$, zero-confinement $2^{S_n}\le3^n$ becomes
\[
k_\ell\;\ge\;\ell\beta\quad\text{at every }\ell,\qquad \beta=1/\alpha=\ln2/\ln3,
\]
so \emph{confinement is exactly a lower bound on ones-density, uniform in the
prefix length}. This is why the critical Sturmian word --- whose ones-count
$\lceil\ell\beta\rceil$ is the least possible at every $\ell$ --- is the extremal
confined word, and why the density theorems below are the right tool for the
non-critical cases.

\medskip
\noindent\textbf{Combinatorial structure: the conjugacy map.} Bernstein and
Lagarias \cite{BernsteinLagarias1996,Bernstein1994} showed that
$x\mapsto v(x)$, the parity vector, is a bijection onto $\{0,1\}^{\N}$ and a
$2$-adic isometry (\lab{cited}); we write $\Ph$ for its inverse. Rationality of
$\Ph(v)$ corresponds to eventual periodicity of $v$ in one direction only; the
converse is Lagarias's Periodicity Conjecture \cite[\S2]{Lagarias1985}
(\lab{cited}, \lab{open}). Under $\Ph$, confinement becomes a density condition:
a zero-confined word has ones-density at least $\beta=1/\alpha=\ln2/\ln3$.

For Sturmian words we follow the conventions of Lothaire
\cite[Ch.~2]{Lothaire2002} (\lab{cited}).

\begin{theorem}[critical Sturmian irrationality]\label{thm:sturmian}
$\Ph(1c_\beta)\notin\Q$, where $1c_\beta$ is the Sturmian word of slope $\beta$
and intercept $0$. The result is effective: $\Ph(1c_\beta)$ equals no $u/v$ in
lowest terms with $\max(|u|,v)\le 2^{301973}$.
\end{theorem}
\lab{proved (paper)} \cite{DeJesusSturmian}. The periodic approximants and their
approximation depths are not new: for every irrational slope
\cite[Thm.~1]{LopezStoll2009} (\lab{cited}) gives a closed $2$-adic series for
$\Ph(1c_\gamma)$ indexed by consecutive pairs of continued-fraction convergents,
whose term exponents $q_{j+1}+q_j-1$ are exactly the depths used; what is new in
\cite{DeJesusSturmian} is the height comparison that turns them into an
irrationality proof. The same constant appears in the author's earlier notes in a
different normalisation, as the $2$-adic Sturmian carry constant
\cite{DeJesusA} and at the Sturmian--Mahler edge of the realizer problem
\cite{DeJesusB}. The word $1c_\beta$ is the
\emph{extremal} zero-confined parity word, its ones-count $\lceil \ell\beta\rceil$
being the smallest compatible with confinement at every length. An extension to
every irrational slope, with $n_0=3$ uniformly and margin
$c(\gamma)=2-\max(1,\gamma\log_2 3)\ge 2-\log_2 3>0$, is
\lab{proved (paper)}, \emph{in preparation}.

\begin{remark}[credit for the general-slope statement]\label{rem:credit}
For irrational $\gamma<\beta$ the irrationality of $\Ph(1c_\gamma)$ is an
immediate consequence of Monks and Yazinski
\cite[Thm.~2.7(b)]{MonksYazinski2004} (\lab{cited}, refereed), a deduction drawn
explicitly by L\'opez and Stoll \cite[p.~6]{LopezStoll2021}. For $\gamma>\beta$
the corresponding statement is \emph{claimed} in the first half of
\cite[Thm.~1]{LopezStoll2021}, resting on that preprint's own archimedean
argument; \textbf{that source is an unrefereed preprint with a single version and
no journal reference, and we treat its half as a claim, not a known theorem,
wherever it appears below.} Only at $\gamma=\beta$ is the statement new.
\end{remark}

\begin{theorem}[density exclusions]\label{thm:density}
Let $v$ be aperiodic. If $\liminf_n k_n(v)/n<\beta$ then $\Ph(v)\notin\Q$
(\lab{cited}, refereed: \cite[Thm.~2.7(b)]{MonksYazinski2004}). The same
conclusion for $\liminf_n k_n(v)/n>\beta$ is \lab{cited} from
\cite[Thm.~1]{LopezStoll2021}, an \textbf{unrefereed} preprint. Consequently the
parity word of a hypothetical zero-confined integer orbit has lower ones-density
\emph{exactly} $\beta$.
\end{theorem}

\begin{proposition}[balanced-square criterion]\label{prop:squares}
Suppose an aperiodic parity word has arbitrarily long prefixes $WW$ whose block
$W$, of length $\ell$ with $k$ ones, is \emph{balanced}. Put
$\theta_W=(k/\ell)\log_2 3$. Then $\Ph(v)\notin\Q$, by
\[
\bigl(2-\max(1,\theta_W)\bigr)\,\ell\ \le\ \log_2 H+O(\log \ell),
\]
$H$ the height of $\Ph(v)$. The balance hypothesis cannot be dropped: for
$W=0^{a}1^{a}$ the ratio $c_W/\bigl(\ell\max(2^{\ell},3^{k})\bigr)$ is unbounded,
and with only the trivial height bound the displayed inequality is vacuous on
confined prefixes.
\end{proposition}
\lab{proved (paper)}. The mechanism is the one Adamczewski and Bugeaud
\cite{AdamczewskiBugeaud2006} exploit in the archimedean setting, resting on the
combinatorics of Allouche, Davison, Queff\'elec and Zamboni \cite{ADQZ2001} and
Berth\'e, Holton and Zamboni \cite{BertheHoltonZamboni2006} (\lab{cited}),
transplanted to the $2$-adic place:
repetitions give good rational approximations, and a height bound forbids them.
The criterion is \emph{one-sided}. It excludes words that are rich in balanced
squares; it says nothing about a word that has none, and a hypothetical confined
word has no reason to have any.

\medskip
\noindent\textbf{Dynamical structure: descent.}

\begin{theorem}[the descent lemmas]\label{thm:descent}
\textup{(L1)} If $m$ is zero-confined for $N$ steps and $m\equiv2\ (\mathrm{mod}\ 3)$,
then $p=(2m-1)/3$ is a positive odd integer with $p<m$, $T(p)=m$, and $p$ is
zero-confined for $N+1$ steps.
\textup{(L2)} The backward $d=1$ chain from $x$ has length exactly $\vt[3](x+1)$,
with $3^{j}\bigl(\mathrm{back}^{j}x+1\bigr)=2^{j}(x+1)$.
\textup{(L4)} $r_{\min}(N,0)\equiv3$ or $7\ (\mathrm{mod}\ 12)$ for every $N\ge1$.
\textup{(L7)} Under a forward step with $d=1$, both $\vt[3](m_k+1)$ and the
round-trip requirement advance by exactly one, so their difference is invariant;
a step with even $d$ destroys the budget.
\end{theorem}
\lab{proved (lean)}, the \texttt{Descent} library. \lab{verified} for (L4) at all
$346$ certified values of $N$, where the residue census is
$3\ (\mathrm{mod}\ 12)$: $218$, $7$: $128$, $11$: $0$.

\begin{proposition}[L4$'$]\label{prop:l4prime}
For $i\ge2$, if $r_i(N)\equiv2\ (\mathrm{mod}\ 3)$ then $(2r_i(N)-1)/3$ is one of
$r_1(N),\dots,r_{i-1}(N)$; and $r_1(N)\not\equiv2\ (\mathrm{mod}\ 3)$, which is
\textup{(L4)}.
\end{proposition}
\lab{proved (paper)}, immediately from (L1). \lab{verified} at every $N\le346$.

\medskip
\noindent The four exclusions, as a map:

\begin{center}\small
\begin{tabular}{@{}llp{0.40\textwidth}@{}}
\toprule
kind of structure & excluded by & what survives it\\
\midrule
algebraic / periodic & Thms.~\ref{thm:cycledrift}, \ref{thm:marker}
& nothing positive: periodic confined points are negative rationals\\
combinatorial (density) & Thm.~\ref{thm:density}
& only lower density \emph{exactly} $\beta$, and the half above $\beta$ rests on
an unrefereed preprint\\
approximation-theoretic & Thm.~\ref{thm:sturmian}, Prop.~\ref{prop:squares}
& words with no long balanced repetition; a hypothetical confined word has no
reason to have one\\
dynamical (descent) & Thm.~\ref{thm:descent}(L7)
& everything: the $3$-adic budget cannot be accumulated, so descent removes no
candidate at all\\
\bottomrule
\end{tabular}
\end{center}

\medskip
\noindent\textbf{The map, read off.} Periodicity is excluded by drift; algebraic
rationality of the conjugacy value is excluded at the critical slope and, granting
one unrefereed half, at every non-critical density; approximation-theoretic
structure --- repetitions --- is excluded when the repeated blocks are balanced;
and dynamical structure in the form of descent is excluded by (L7). Every one of
these exclusions removes a \emph{family} of candidates. None of them removes a
candidate that has no structure at all, and Section~\ref{sec:needed} argues that
this is not an accident of the methods used.

\section{Closed routes}\label{sec:closed}

Each row is a route that was pursued and closed, with the reason it closes. The
purpose of the table is to prevent re-derivation.

\begin{center}\small
\begin{tabular}{@{}p{0.27\textwidth}p{0.55\textwidth}l@{}}
\toprule
route & why it fails & label\\
\midrule
descent by inheritance and well-ordering
& The round-trip criterion is \emph{provably equivalent} to (DE); and (L7) shows
the $3$-adic budget $\vt[3](m+1)$ cannot be accumulated by forward motion, since a
$d=1$ step raises budget and requirement together while an even $d$ destroys the
budget. Observed round trips are no more frequent than a geometric baseline.
& \lab{proved (paper)}\\[2pt]
transport budgets
& The budget identity collapses to the aggregate identity; the bookkeeping is a
restatement.
& \lab{proved (paper)}\\[2pt]
time-axis and word-level arguments
& The $2$-adic fixed point $-1$ is zero-confined forever, so no argument at the
level of words and drift alone can exclude confinement; and uncountably many
confined words satisfy every word-level shadow while being unrealizable by any
integer.
& \lab{proved (paper)}\\[2pt]
fixed-depth Fourier and triangle methods
& The implication from a dangerous-window sparsity hypothesis to low-frequency
decay is machine-checked, but the hypothesis itself is open and the frequency
$\lambda=1$ is unavoidable, so no frequency averaging is available where it is
needed.
& \lab{proved (lean)} for the implication; hypothesis \lab{open}\\[2pt]
parity-type counting on the cycle side
& Digit counting for powers of $2$ in base $3$ gives $o(L)$ information only, far
short of what cycle exclusion needs.
& \lab{proved (paper)}\\[2pt]
record-holder anatomy
& The extremes are generic: every candidate pattern reduces to (L1), (L2), (L4),
(L4$'$), to the last-maximum structure, or to a selection effect. The record
holders are not Sturmian-adjacent, sharing only one to three valuations with the
best shift of the critical Sturmian word over depths up to $236$.
& \lab{proved (paper)}, \lab{verified}\\
\bottomrule
\end{tabular}
\end{center}

\medskip
\noindent A note on the fourth row, because it is the one most often expected to
work. The analytic side of the programme follows the architecture of Tao
\cite{Tao2022} (\lab{cited}), whose decay estimate is uniform in the frequency and
polynomial in strength. The law that arises here is \emph{critical and confined},
and uniform decay fails for it numerically at small integer frequencies. Tao's
result therefore does not transfer directly, and no part of it is used as an
input. That is a statement about the present architecture, not an impossibility
theorem.

\section{What is needed}\label{sec:needed}

\subsection*{(a) Three conditions, and the controls a candidate must fail}

\begin{criterion}[the three conditions]\label{crit:three}
A tool capable of deciding \textup{(DE)} must
\begin{enumerate}
\item[(i)] use $2$-adic integrality at unboundedly many depths;
\item[(ii)] use the sign, or the archimedean size, of the integer; and
\item[(iii)] couple (i) and (ii) along a \emph{single} orbit.
\end{enumerate}
\end{criterion}
\lab{heuristic}: this is a checklist distilled from the closed routes of
Section~\ref{sec:closed}, not a theorem. Its use is negative and it is cheap to
apply. The statistics of Section~\ref{sec:statistics} fail (ii) and (iii), having
no specific integer in them. The irrationality results of
Section~\ref{sec:excluded} satisfy (i) but fail (ii) and (iii): they use only that
an integer $M_n$ is nonzero, never its sign, and they concern a single fixed word
rather than an orbit of integers. The descent lemmas satisfy (ii) and (iii) but
fail (i): their budget is $3$-adic and (L7) proves it cannot accumulate. The
reduction of Theorem~\ref{thm:reduction} satisfies all three and yields nothing,
because it is an equivalence.

\medskip
\noindent\textbf{The checklist applied, once.} As an illustration, take the most
natural proposal a reader is likely to make: \emph{a divergent orbit has drift
$R_n\le0$ forever, hence its valuation mean is at most $\alpha$; combine this with
a counting bound on how many integers below $X$ can have such a mean, and conclude
that no integer does}. This satisfies (i), since confinement at every depth is a
$2$-adic integrality condition at unbounded depth. It fails (ii): no sign and no
size of any particular integer enters, only a count. And it fails (iii): the count
is over a population, not along one orbit. The proposal is
Section~\ref{sec:statistics} with a different constant, and
Section~\ref{sec:statistics} shows the population is exactly as large as random
placement predicts --- so a counting bound of this kind cannot be improved to a
contradiction, because there is no discrepancy to exploit.

\medskip
\noindent\textbf{The controls.} A candidate tool must \emph{fail} on each of the
following, and checking this is the first thing to do with any proposal.
\begin{itemize}
\item $-1$ under $3x+1$: a zero-confined fixed point. Any argument that excludes
positive confinement without using positivity will exclude it too, and is wrong.
\item $\{-5,-7\}$ under $3x+1$: a zero-confined $2$-cycle. Note that confinement
is a property of a \emph{point on} the cycle, not of the cycle: $-5$ is confined
forever while $-7$, on the same cycle, fails at the first step. Any argument
insensitive to the starting point is wrong.
\item $+1$ under $3x-1$: a zero-confined \emph{positive} fixed point. See
Theorem~\ref{thm:signobstruction}.
\item random confined words: their least realizers sit at the class modulus
$2^{S_N}$, and an argument that excludes them excludes almost everything and has
proved too much.
\end{itemize}

\subsection*{(b) The sign obstruction, as a theorem}

\begin{theorem}[sign obstruction]\label{thm:signobstruction}
No argument that uses only the aggregate identity, the confinement condition and
the sign of the denominator can separate the positive integers from the rest of
$K$.
\end{theorem}

\begin{proof}[Reason]
The map $m\mapsto -m$ is an exact conjugacy from the accelerated $3x-1$ map on the
positive integers to the accelerated $3x+1$ map on the negative integers,
valuation by valuation: $\vt(3(-m)+1)=\vt(1-3m)=\vt(3m-1)$ and $T(-m)=-T'(m)$,
where $T'(m)=(3m-1)/2^{\vt(3m-1)}$. The two systems therefore have identical
valuation words, identical $S_k$ and identical confinement. Now run the
computation of Theorem~\ref{thm:marker} in the $3x-1$ system. The one-step
identity is $2^{a_k}m_{k+1}=3m_k-1$, so the carry moves to the other side of the
periodicity equation, $m(3^{L}-2^{S})=C'_L$ with the \emph{same} positive
$C'_L=\sum_{i<L}3^{\,L-1-i}2^{\,S_i}$, and confinement gives $3^{L}-2^{S}>0$,
whence $m>0$. Indeed $+1$ is a zero-confined positive fixed point of $3x-1$, and
its record ladder is degenerate: $r_1(N)=1$, $r_2(N)=5$, $r_3(N)=17$ for every
$N$ up to the computed cap, these being the images of $-1$, $-5$, $-17$. The two
systems agree on the aggregate identity, on the confinement condition and on the
$S_k$ combinatorics, and disagree on the answer. Whatever separates them is
something the listed ingredients do not use.
\end{proof}
\lab{proved (paper)}; the conjugacy \lab{verified} on every odd $m<4000$ to depth
$300$, identical words and depths. Compare the descent audit, which found the same
sensitivity from the backward side: the backward fixed point of $3x+1$ is $-1$ and
of $3x-1$ is $+1$, so the strict decrease $p<m$ of (L1) is unconditional in the
first case and fails at exactly one point in the second.

\subsection*{(c) Border versus height}

\begin{formulation}[the price of agreement]\label{form:border}
Every exclusion in Section~\ref{sec:excluded} has one shape. If a $2$-adic target
$x$ agrees with a rational of height $H$ to depth $R$ --- the \emph{border} --- then
\[
R\ \le\ c\cdot\log_2 H+O(\log R)
\]
with $c$ depending on the structure of the approximant. \textup{(DE)} asserts that
a \emph{finite} height cannot buy an \emph{infinite} border inside $K$: an integer
$m$ is a rational of height $|m|$, and being zero-confined forever means agreeing
with the confinement corridor at every depth.
\end{formulation}
\lab{heuristic} as a programme statement. Two remarks on what cannot replace
height. First, descriptive complexity does not: the Liouville step consumes
$\log_2 H$ and nothing weaker, because the integer $M_n$ it produces is bounded
below by $2^{R}$ and above by $H$ times a polynomial, and no measure of
description length enters either bound. It is worth saying what a height bound would buy, since it fixes the exchange
rate. By \cite[Thm.~5.2]{DeJesusB} (\lab{proved (paper)}), a finite effective
irrationality measure $\mu$ for $\Ph(1c_\beta)$ over the integers would give
$r(D_N)\ge(2C)^{1/\mu}2^{S_N/\mu}$ --- an \emph{exponential} floor on the least
integer realising the length-$N$ prefix, far beyond anything the linear scale of
the present estimates provides. That is the shape of a tool satisfying
Criterion~\ref{crit:three}(i) and (ii) at once, and it is open in exactly the form
that would be useful. Second, the word side alone does not: there
are uncountably many confined words satisfying every word-level constraint that
has been extracted, and almost all of them are realized by no integer at all. The
height is the only quantity in the problem that distinguishes an integer from a
$2$-adic limit of integers.

\subsection*{(d) The $\sqrt n$ formulation}

\begin{formulation}[the one-sided walk]\label{form:walk}
By Theorem~\ref{thm:density}, a divergent orbit of a rational $2$-adic integer has
parity word of lower density exactly $\beta$ --- refereed below $\beta$, and
resting on an unrefereed preprint above it. At that density the drift $R_n$ is a
zero-mean lattice walk, and zero-confinement says it is \emph{one-sided forever},
$R_n\le0$ for all $n$. For a genuine random walk this has probability zero, and is
excluded quantitatively by the law of the iterated logarithm.
\end{formulation}
\lab{heuristic}, and we state the caveat inside the formulation rather than after
it: \textbf{this is a reformulation of \textup{(DE)}, not progress toward it}. The
walk is deterministic, generated by the orbit itself; no independence, no mixing
and no martingale structure is available, and the confined words are precisely the
measure-zero set the law of the iterated logarithm describes. What the
reformulation does contribute is a scale: it says the obstruction is of size
$\sqrt n$ in the drift, which is what Theorem~\ref{thm:scales}'s $N^{-3/2}$ is
measuring on the ensemble side.

\subsection*{(e) The only integer-specific features}

\begin{observation}\label{obs:integer}
The only features an integer possesses that a general element of $\Zt$ does not
are its finite bit length --- equivalently, a terminating $0^{\infty}$ tail in its
binary expansion --- and its archimedean size. Every attempt recorded in this
programme to read either through the dynamics has collapsed into a restatement of
\textup{(DE)}.
\end{observation}
\lab{open}. To be concrete about what is missing, here is the shape such a law
would have to take. Let $m<2^{B}$ be odd, so that its binary expansion terminates
after $B$ digits. The parity digits $v_0,\dots,v_{n-1}$ of $m$ are determined by
$m \bmod 2^{n}$, and conversely determine it; the finiteness of $m$ is the
statement that all digits beyond $B$ vanish, which the parity coordinates do not
see directly. A usable law would bound, as a function of $B$ and $n$, how much of
the parity prefix is \emph{forced} by $m<2^{B}$ --- equivalently, how the
terminating tail constrains the orbit's valuation word at depth $n\gg B$. Every
statement of that kind available here is either vacuous (the constraint is
automatic) or equivalent to (DE) itself.
Form (d) of Proposition~\ref{prop:equivalent} is exactly the statement
that the zero tail is what is missing, and it is equivalent to (DE) rather than a
route to it; the height formulation (c) is the size, and it is likewise
equivalent. A genuinely new ingredient would have to be a \emph{quantitative law}
for how the terminating tail propagates into the parity digits --- how many
parity digits of $m$ are constrained by the fact that $m<2^{B}$, as a function of
$B$ and of the depth. Nothing of this kind is known here, and the programme
contains no candidate for it.

\subsection*{(f) The signed ladder of targets}

\begin{proposition}[the ladder]\label{prop:ladder}
Let \textup{(PosPC)} be the assertion that $K$ contains no positive rational.
Then
\[
\text{Lagarias's Periodicity Conjecture}\ \Longrightarrow\ \textup{(PosPC)}
\ \Longrightarrow\ \textup{(DE)},
\]
and the second implication is strict as far as anything known. Moreover
\textup{(PosPC)} holds if and only if, for every odd $v\ge1$, no positive odd
integer has a zero-confined orbit under $n\mapsto(3n+v)/2^{\vt(3n+v)}$.
\end{proposition}
\lab{proved (paper)}. The first implication uses Theorem~\ref{thm:marker}: under
the Periodicity Conjecture every rational $2$-adic integer is eventually cyclic,
and every eventually periodic point of $K$ is negative. The equivalence at the
end is an exact conjugacy, valuation by valuation --- scaling by an odd $v$
preserves every $a_j$ on the nose, which matters because confinement is a
condition at \emph{every} index. So (PosPC) is (DE) asserted simultaneously for
the whole family $3x+v$, and it is the natural signed target sitting strictly
between (DE) and a classical conjecture.

\subsection*{(g) The cycle side}

\begin{observation}\label{obs:cycle}
Cycle exclusion is an anti-concentration statement modulo $2^{S}-3^{L}$: a
nontrivial cycle of length $L$ and total $S$ requires the carry $C_L$ to be
divisible by, and comparable in size to, that modulus. Fixed-prime digit counting
does not reach it. The closest available machinery is the $p$-adic
Subspace/Ridout family \cite{Ridout1958,Mahler1957} (\lab{cited}), which gives
$o(L)$-type information and is \emph{ineffective}: the available proofs do not
determine the exceptional subspaces, so no rate and no computable threshold
follow.
\end{observation}
\lab{open}. Eliahou's cycle-length bounds \cite{Eliahou1993} (\lab{cited}) are
what is actually available on this side, and they are of a different nature: they
constrain $L$ and $S$ through continued-fraction approximations to $\log_2 3$
rather than through anti-concentration.

\section{Open problems}\label{sec:open}

\begin{enumerate}
\item \textbf{(DE)}: no positive odd integer is zero-confined. \lab{open}.
Equivalently, by Theorem~\ref{thm:reduction}, no positive odd integer has a
divergent accelerated orbit.
\item \textbf{The integer irrationality measure.} Is there $\mu<\infty$ and $C>0$
with $|\Ph(1c_\beta)-y|_2\ge C|y|^{-\mu}$ for every \emph{integer} $y\ge1$?
\lab{open}. The rational-approximation analogue is false unless $\log_3 2$ is
badly approximable --- the depth law forces any such exponent to exceed
$1+q_{n+1}/q_n$, and the $331$st partial quotient of $\log_3 2$ is $2436$ --- but
the integer form, which is the one that converts into an exponential lower bound
on least realizers \cite[Thm.~5.2]{DeJesusB}, is untouched by that obstruction.
\item \textbf{Sharp EOC}, Conjecture~\ref{conj:sharp}. \lab{open}.
\item \textbf{(PosPC)}: $K$ contains no positive rational.
\lab{open}; strictly between (DE) and Lagarias's Periodicity Conjecture by
Proposition~\ref{prop:ladder}.
\item \textbf{Sturmian words at general intercept.} Is $\Ph(1c_{\gamma,\rho})$
irrational for every irrational $\gamma$ and every $\rho$? \lab{open} for
$\rho\ne0$. The obstruction is precise: the one-positions are no longer
$\lfloor j\alpha_\gamma\rfloor$, and the two exact identities on which the depth
law turns lose their right-hand sides.
\item \textbf{Transcendence of $\Ph(1c_\beta)$.} \lab{heuristic}. The depth law
supplies rational approximations of exponent $1+q_{n+1}/q_n$, which exceeds the
Roth threshold of $2$ along the odd convergents; whether the applicable form of
the $p$-adic Roth theorem \cite{Ridout1958,KalitzinMurru2026} can consume them,
and whether the statement is already in the literature, has not been checked. It
is recorded here as a direction, not a result.
\item \textbf{Cycle-side anti-concentration} modulo $2^{S}-3^{L}$, in the sense of
Observation~\ref{obs:cycle}. \lab{open}.
\item \textbf{Dangerous-window sparsity} for the analytic side: the hypothesis
that implies low-frequency decay, for every shift and some logarithmic window
length. \lab{open}.
\end{enumerate}

\noindent None of these, and nothing else in this note, bears on the Collatz
conjecture.

\section{Formal and computational resources}\label{sec:resources}

\noindent\textbf{Repositories.} \texttt{eoc-divergence} \cite{EOCDivergenceRepo}
holds three independent Lean~4 libraries and four audits;
\texttt{eoc-lean-verification} \cite{EOCLeanRepo} holds the analytic side.

\medskip
\noindent\textbf{Lean libraries} (Lean \texttt{v4.34.0}, Mathlib
\texttt{v4.34.0}; no \texttt{sorry}, \texttt{axiom}, \texttt{opaque},
\texttt{native\_decide} or \texttt{implemented\_by}; every headline theorem
audited to depend on \texttt{propext}, \texttt{Classical.choice},
\texttt{Quot.sound} only, enforced in continuous integration).
\texttt{Divergence} ($\approx1650$ lines) proves Theorems~\ref{thm:reduction},
\ref{thm:sparsity}, \ref{thm:summable}, \ref{thm:cycledrift}.
\texttt{Occupation} ($\approx1430$ lines) proves Theorem~\ref{thm:rate} and the
cycle lemma; it is not imported by \texttt{Divergence}.
\texttt{Descent} ($\approx430$ lines) proves Theorem~\ref{thm:descent}.

\medskip
\noindent\textbf{Audits.} \texttt{audits/descent} (verdict \textsc{stop});
\texttt{audits/sturmian\_irrationality}, including the general-slope phase;
\texttt{audits/record\_holder\_anatomy}, which produced
Propositions~\ref{prop:density}, \ref{prop:clustering}, \ref{prop:l4prime} and
Theorem~\ref{thm:signobstruction}. Each carries its own report with every claim
labelled and every computation scripted.

\medskip
\noindent\textbf{Reproducibility.} All arithmetic reported here is exact: integer
or rational, never floating-point, in every step that reaches a decision. The
record-holder scan is a C program testing $S_k\le A[k]$ with
$A[k]=\mathrm{bitlength}(3^k)-1$ precomputed exactly and orbit values carried in
\texttt{unsigned \_\_int128} behind an overflow guard; it covers every odd
$m\equiv3\ (\mathrm{mod}\ 4)$ below $10^{12}$ in about $37$ minutes on eight
cores, and was cross-checked three ways --- against previously published values
for $N\le200$, against an independent arbitrary-precision reference on random
windows, and against a residue-class reconstruction of each record holder from its
own valuation word. The density $p_N(0)$ is computed by the transfer recursion of
Proposition~\ref{prop:density} and cross-checked against direct orbit counting.

\section*{Declaration of AI assistance}

AI assistance (Anthropic's Claude, through Claude Code) was used throughout the
programme for adversarial auditing of proposed arguments, literature search and
verification of bibliographic data, for writing and running the exact computations
reported here, and for drafting. In particular, the record-holder audit of
Section~\ref{sec:statistics}, the sign-obstruction analysis of
Section~\ref{sec:needed} and the present note were drafted in that setting under
pre-registered protocols with stated stop rules; several claims were corrected or
withdrawn as a result, and those corrections are recorded in the audit reports
rather than silently applied. All mathematical claims, their labels, their
framing, and any errors are the author's.

\bibliographystyle{amsplain}
\bibliography{refs}

\end{document}
